% !TeX root = ../sections/02_exercises.tex
\subsection{2-B-2.}
\begin{exercise}
	有界な集合 \(A,B\subset\mathbb{R}\) に対して，
	\[
	\sup(A\cup B)
	=
	\max\{\sup A,\sup B\}
	\]
	および
	\[
	\inf(A\cup B)
	=
	\min\{\inf A,\inf B\}
	\]
	が成立することを証明せよ。
\end{exercise}

\begin{background}
	この問題では，和集合の上限と下限を扱う。
	
	集合 \(A\cup B\) の元は，
	\(A\) の元または \(B\) の元である。
	したがって，\(A\cup B\) の上界になるためには，
	\(A\) の上界でもあり，かつ \(B\) の上界でもある必要がある。
	
	よって，\(A\cup B\) の最小上界は，
	\[
	\sup A
	\quad\text{と}\quad
	\sup B
	\]
	のうち大きい方になる。
	
	同様に，\(A\cup B\) の下界になるためには，
	\(A\) の下界でもあり，かつ \(B\) の下界でもある必要がある。
	したがって，最大下界は
	\[
	\inf A
	\quad\text{と}\quad
	\inf B
	\]
	のうち小さい方になる。
\end{background}

\begin{strategy}
	まず，上限について示す。
	\[
	\alpha=\sup A,
	\qquad
	\beta=\sup B,
	\qquad
	M=\max\{\alpha,\beta\}
	\]
	とおく。
	
	任意に \(x\in A\cup B\) を取る。
	このとき \(x\in A\) または \(x\in B\) である。
	もし \(x\in A\) ならば
	\[
	x\leq\alpha\leq M
	\]
	であり，もし \(x\in B\) ならば
	\[
	x\leq\beta\leq M
	\]
	である。
	したがって，\(M\) は \(A\cup B\) の上界である。
	
	次に，\(\gamma\) を \(A\cup B\) の任意の上界とする。
	すると，\(A\subset A\cup B\) かつ \(B\subset A\cup B\) なので，
	\(\gamma\) は \(A\) の上界でもあり，\(B\) の上界でもある。
	したがって
	\[
	\alpha\leq\gamma,
	\qquad
	\beta\leq\gamma
	\]
	となる。
	よって
	\[
	M=\max\{\alpha,\beta\}\leq\gamma
	\]
	である。
	したがって
	\[
	\sup(A\cup B)=M
	\]
	が従う。
	
	下限についても同様に示す。
	\[
	i=\inf A,
	\qquad
	j=\inf B,
	\qquad
	m=\min\{i,j\}
	\]
	とおく。
	任意の \(x\in A\cup B\) に対して，
	\(x\in A\) ならば
	\[
	m\leq i\leq x
	\]
	であり，\(x\in B\) ならば
	\[
	m\leq j\leq x
	\]
	である。
	したがって \(m\) は \(A\cup B\) の下界である。
	
	また，\(\delta\) を \(A\cup B\) の任意の下界とすると，
	\(\delta\) は \(A\) と \(B\) の下界でもある。
	よって
	\[
	\delta\leq i,
	\qquad
	\delta\leq j
	\]
	である。
	したがって
	\[
	\delta\leq\min\{i,j\}=m
	\]
	となる。
	よって
	\[
	\inf(A\cup B)=\min\{\inf A,\inf B\}
	\]
	である。
\end{strategy}